%! Exponential Growth and Decay — Unit 6, Lesson 6.1 — Lesson Plan
\documentclass[10pt]{article}
\usepackage{algebra2-article}
\usepackage{algebra2-boxes}
\usepackage{graphicx}   % -article does NOT load graphicx; needed for thumbnails/screenshots
\graphicspath{{images/}}
\newcommand{\TallMath}[1]{$\displaystyle #1\rule[-1.4em]{0pt}{3.2em}$}
\newcommand{\CourseName}{Algebra 2: Shepherd}
\newcommand{\SchoolYear}{2026--2027}
\newcommand{\MeetingLength}{55 minutes}
\newcommand{\UnitNumberName}{Unit 6: Exponential Functions \quad}
\newcommand{\LessonNumberName}{Lesson 6.1: Exponential Growth and Decay}

\begin{document}
\begin{center}
    {\Large\bfseries \CourseName: \SchoolYear \\
    {\small\bfseries \UnitNumberName \LessonNumberName}}
\end{center}

\begin{tcolorbox}[colback=forestbg, colframe=forest, boxrule=0.9pt, arc=2mm,
    left=2mm, right=2mm, top=1.5mm, bottom=1.5mm]
    \textbf{Primary Objective:} Students can \emph{build} an exponential model $y=ab^{x}$ from any of
    the three sources it usually arrives in --- a \emph{table} (where $a$ is the $x=0$ output and $b$
    is the constant ratio), \emph{two points} (where dividing the two equations cancels $a$ and
    leaves $b$ raised to the gap between the inputs), and a \emph{written description} (where a
    percent rate $r$ must be translated into a factor). The lesson turns on that translation:
    $b=1+r$ for growth and $b=1-r$ for decay, so ``up $8\%$'' gives $b=1.08$ and ``down $15\%$''
    gives $b=0.85$ --- the base is the fraction you \emph{end up with}, never the percent that
    changed. Students then \emph{evaluate} the model they built and \emph{interpret} the result in
    context, including judging how far out the model stays believable.
    \\[2pt]
    \textbf{Standards (2023 VA SOL):} \textbf{A2.F.2f} --- evaluate a function at a value in its
    domain and \emph{interpret the result in context} --- is the assessed standard of the day; every
    model students build is then used and read back into its story (the town, the car, the bacteria
    culture, the phone, the two towns). \textbf{A2.F.2a} and \textbf{A2.F.2h} are carried forward
    from Lesson~6.0 whenever a range or a horizontal asymptote is checked as part of validating a
    model (Homework 3, Activity Tier~A and Tier~E). \textbf{A2.F.1e} (compare and contrast graphs,
    tables, and equations) is the structural spine of the whole lesson --- the same function is
    built three different ways on purpose. \textbf{A2.F.2b} drives the linear-versus-exponential
    contrast in the Hook and in Activity Tier~E Part~3. This lesson is also the \emph{prerequisite}
    for \textbf{A2.ST.2d/e} in Lesson~6.5: a regression is worthless if students cannot say what $a$
    and $b$ mean once technology hands them over. \emph{Note:} Unit~6 has no \texttt{A2.EO} or
    \texttt{A2.EI} standard --- the whole unit rests on \texttt{A2.F} and \texttt{A2.ST.2}.
\end{tcolorbox}

\renewcommand{\arraystretch}{1.4}
\begin{skillbox}[Priority Ideas \& Skills]{goldbox}
\begin{tabularx}{\linewidth}{@{}X|X@{}}
\textbf{Priority Ideas \& Skills} & \textbf{Key Understandings (the ``why'')} \\[2pt]
\hline
\vspace{-6pt}
\begin{itemize}[leftmargin=1.1em, nosep, after=\vspace{2pt}]
  \item Write $y=ab^{x}$ from a \textbf{table}: $a$ from the $x=0$ column, $b$ from the constant ratio.
  \item Translate a percent into a factor: $b=1+r$ for \textbf{growth}, $b=1-r$ for \textbf{decay}.
  \item Write $y=ab^{x}$ from a \textbf{story}, naming what $a$ and $b$ mean in that context.
  \item Write $y=ab^{x}$ from \textbf{two points}, including when neither is the $y$-intercept.
  \item \textbf{Evaluate} the model and \textbf{interpret} the output as a sentence about the situation.
  \item Run the reverse translation: read a percent rate back off a finished model.
  \item \textbf{Validate} a model --- does $x=0$ return the starting amount, and is $b$ on the correct side of $1$?
  \item Diagnose a wrong model by saying what it \emph{actually} describes.
\end{itemize}
&
\vspace{-6pt}
\begin{itemize}[leftmargin=1.1em, nosep, after=\vspace{2pt}]
  \item \emph{The base is the fraction you end up with, not the percent that changed.} One sentence; every error today violates it.
  \item Up $8\%$ leaves you with $108\%$. Down $15\%$ leaves you with $85\%$. That is where $1+r$ and $1-r$ come from --- they are not rules to memorize.
  \item Building is reading run backwards. Lesson 6.0 pulled $a$ and $b$ out of a graph; today the same two numbers get pushed in.
  \item A percent applies to the \emph{current} amount, not the original one. That is the entire difference between exponential and linear growth.
  \item Two points determine the curve because dividing kills $a$ --- and the leftover exponent is the \emph{gap} between the inputs, not one of them.
  \item A model is a claim about the world, so it can be checked \emph{and} disbelieved. Both belong in the answer.
  \item A wrong base is not just wrong --- it is right about a \emph{different} situation. Naming that situation is how you find the error.
\end{itemize}
\end{tabularx}
\end{skillbox}

\begin{skillbox}[{Vocabulary, Concepts, \& Theorems}]{sky}
\renewcommand{\arraystretch}{1.3}
\begin{tabularx}{\linewidth}{@{}l X@{}}
\textbf{Initial value $a$} & The amount the model starts with: the output at $x=0$, and the $y$-intercept of the graph. \\
\textbf{Growth / decay factor $b$} & What you multiply by for each one-unit step. $b>1$ grows; $0<b<1$ decays. \\
\textbf{Growth / decay rate $r$} & The percent change per period, written as a decimal: $8\%$ becomes $r=0.08$. \\
\textbf{Exponential growth model} & $y=a(1+r)^{t}$ --- increasing by the same \emph{percent} each period, so $b=1+r$. \\
\textbf{Exponential decay model} & $y=a(1-r)^{t}$ --- decreasing by the same \emph{percent} each period, so $b=1-r$ is the fraction that \emph{remains}. \\
\textbf{Percent-to-factor rule} & Add or subtract $r$ from $1$; never use $r$ itself as the base. $r=0.15$ decay gives $b=0.85$, not $0.15$ and not $1.15$. \\
\textbf{Constant ratio (from 6.0)} & Dividing consecutive outputs gives the same number every time --- and that number \emph{is} $b$. \\
\textbf{Building from two points} & Substituting both points and dividing cancels $a$, leaving $b$ raised to the difference of the inputs. \\
\textbf{Contextual domain} & The inputs the story actually allows, which is usually narrower than the algebraic domain of all reals. \\
\end{tabularx}
\end{skillbox}

\begin{fixedskillbox}[Activate Prior Knowledge \& Spiral Review]{forestbg}
    The Warm-Up rehearses exactly what today needs, in the order the lesson uses it. \textbf{(1)}
    A \$$50$ jacket goes up $8\%$ and then (starting over) down $15\%$, worked \emph{twice} --- once
    by finding the change and adding or subtracting it, once as a single multiplication --- and
    finished with ``the new price is \rule{0.7cm}{0.4pt}\% of the old price.'' That last line is the
    whole lesson: students who can say ``$108\%$'' will write $b=1.08$, and students who only see
    ``$8\%$'' will write $b=0.08$. \textbf{(2)} A two-row recall check from Lesson~6.0 --- pull $a$,
    $b$, growth/decay, and the $y$-intercept off $y=5(1.2)^{x}$ and $y=80(0.6)^{x}$ --- because
    today is that skill run backwards, and anyone who cannot read a model cannot write one.
    \textbf{(3)} A table with constant ratio $3$ and initial value $7$, which \emph{is} the
    build-from-a-table method before it has a name. Debrief item 1 at the board and put both
    $1.08$ and $0.85$ up where they will stay all period.
\end{fixedskillbox}

\begin{skillbox}[Hook]{forestbg}
    Two towns, same population today. Put them on the board side by side: \textbf{Riverton --- $8{,}000$
    people, ``we expect to gain $300$ people every year.''} \textbf{Kellsboro --- $8{,}000$ people,
    ``we expect to grow by $4\%$ every year.''} Ask which is linear and which is exponential, and
    refuse arithmetic as an answer: the giveaway is \emph{add} versus \emph{multiply}, which is
    Lesson~6.0's fingerprint test in words. Then ask the question the whole lesson hangs on:
    \textbf{``Kellsboro grows $4\%$ this year. Next year it grows $4\%$ again --- of \emph{what}?''}
    Take answers until someone says \emph{of the new, bigger number}. That is why the model
    multiplies. Now write the trap on the board deliberately: \emph{Kellsboro:} $P(t)=8000(0.04)^{t}$
    --- and ask the class to check it at $t=1$. It gives $320$ people. Let the absurdity land, then
    reuse the Warm-Up's phrasing --- growing $4\%$ leaves you with $104\%$ of the town --- and
    rewrite it as $8000(1.04)^{t}$. Do \emph{not}
    generalize to $1+r$ yet --- Section 2 of the notes earns that.
\end{skillbox}

\begin{skillbox}[Lesson]{forestbg}
\begin{multicols}{2}
\textbf{1. Two numbers, and a table that hands you both.} Open by naming the job: building a model
means finding $a$ and $b$, nothing else. Work the table $3,12,48,192,768$ --- $a$ is the $x=0$
output, $b$ is the constant ratio $4$ --- write $y=3\cdot4^{x}$, and then \emph{check it} against
the $x=4$ row students did not use. Make the check a habit today; it is what makes the rest of the
lesson self-correcting. The ``your turn'' table ($800,200,50,12.5$) is where the upside-down
division error appears; catch it now, because a student who divides earlier-by-later turns every
decay model into growth.

\textbf{2. Percent $\to$ factor.} This is the section the lesson is named for. Derive $b=1+r$ and
$b=1-r$ from the Warm-Up rather than announcing them: after an $8\%$ raise you have $108\%$ of what
you had, and $108\%$ \emph{is} $1.08$. Fill the six-row translation table together. Row 5 (``$100\%$
of it remains'') is the illegal base $b=1$ and pays off Lesson~6.0's restriction with a reason
students can feel --- nothing changes, so it is a horizontal line. Row 6 (``doubles'') is $r=100\%$,
which surprises people. Then spend real time on the red error box: $b=0.08$ for ``up $8\%$'' is a
$92\%$ \emph{loss}, and $b=1.15$ for ``down $15\%$'' is a $15\%$ \emph{gain}. Install the two-second
check out loud --- \emph{is $b$ on the correct side of $1$?}

\columnbreak

\textbf{3. Story to model, and back.} Three moves, every time: pull $a$, translate to $b$, write
$y=ab^{x}$. Work the growth case (town of $12{,}000$ at $3\%$, predict $P(10)\approx16{,}127$) and
the decay case (car at \$$24{,}000$ losing $12\%$, $V(5)\approx\$12{,}666$) in full, and require a
\emph{sentence} after each evaluation --- that sentence is the \textbf{A2.F.2f} evidence, and it is
the part students skip. The decay example's last question --- \emph{why $0.88$ and not $0.12$?} ---
is the cold-call of the day. Close with the three-step check-your-model routine.

\textbf{4. Two points.} The hardest page. Case 1 has the $y$-intercept given, so $a$ is free and
$40=5b^{3}$ leaves a cube root --- Unit 5 machinery, and say so. Case 2 uses the graph through
$(-1,1)$ and $(2,8)$: write both points into $y=ab^{x}$, divide, and watch $a$ cancel. Two things
get dropped here and both need saying out loud: the leftover exponent is the \emph{gap} between the
$x$-values ($2-(-1)=3$), and $b$ still has to go back into one of the points to recover $a$. Finish
by checking the model against the graph's $y$-intercept --- the same validation habit from
Section 1. Preview Lesson~6.2.
\end{multicols}
\end{skillbox}

\begin{skillbox}[Active Monitoring]{redbox}
Circulate during the notes practice and Tier work. Watch for and correct:
\begin{itemize}[leftmargin=1.2em, nosep]
  \item using the \textbf{percent as the base} --- $b=0.08$ for ``up $8\%$,'' $b=0.12$ for ``loses $12\%$''; this is \emph{the} signature error of the lesson;
  \item the \textbf{sign} error --- $b=1.15$ for a $15\%$ \emph{decrease}, or $b=0.85$ for an increase; adding when the story subtracts;
  \item dividing the table \textbf{upside down} (earlier output over later), which flips decay into growth;
  \item forgetting that $a$ is the output at $x=0$ specifically --- grabbing the first number in a table whose $x$-values do not start at $0$;
  \item in two-point problems, using one of the $x$-values as the exponent instead of the \emph{gap} between them, or finding $b$ and never going back for $a$;
  \item taking a square root when the gap was $3$ (Homework 3: $b^{3}=\frac18$ gives $\frac12$, not $\frac13$);
  \item reporting a bare number with no sentence --- ``$2940.10$'' is not an interpretation;
  \item modeling a constant \emph{amount} of change with an exponential, or a constant \emph{percent} with a linear model (Activity Tier E).
\end{itemize}
Cold-call prompts: ``After the change, what percent of the original do you have?'' ``Is your $b$ bigger
or smaller than $1$ --- and should it be?'' ``What does your model give at $x=0$? Is that the number
in the story?'' ``Your classmate's model is wrong --- but what situation would it be \emph{right}
for?'' ``The exponent came out $3$. Where did the $3$ come from?''
\end{skillbox}

\begin{skillbox}[Group Work \& Differentiation]{redbox}
\textbf{Building the Model} activity (tiered; mirrors the notes).
\begin{multicols}{3}
\textbf{Tier R --- Remediate.} A four-row percent-to-factor table whose first two rows are the
\emph{same} percent in opposite directions ($6\%$ up, $6\%$ down) with a required explanation of why
the bases differ. Then one model built end to end --- $60$ ants growing $25\%$ a week --- with the
validation steps written in as their own lettered parts: what does the model give at $w=0$, is
$N(1)$ really $25\%$ more than $60$, and what would $60(0.25)^{w}$ describe instead.

\columnbreak
\textbf{Tier A --- Approaching.} The same function built from three different sources: a table
($9,27,81,243$), a pre-drawn graph of $9\left(\frac13\right)^{x}$ with three labeled points, and a
savings-account story at $4\%$. Then the \textbf{Justify}: a deer herd declining $20\%$ modeled three
ways by three students ($1.2$, $0.2$, $0.8$), where identifying Rosa is the easy half and diagnosing
Maya's \emph{sign} error against Devon's \emph{percent-as-base} error is the real assessment ---
two different misconceptions needing two different re-teaches.

\columnbreak
\textbf{Tier E --- Extension.} Three parts. \emph{Two points, neither at $x=0$:} $(2,45)$ with
$(5,1215)$, then $(1,60)$ with $(4,7.5)$ where $b$ comes out a fraction. \emph{Reverse translation:}
read the initial value and percent rate back off four finished models, including $75(2)^{t}$ as
$100\%$. \emph{The two leaks:} a $240$-gallon tank losing $60$ gallons a day versus $25\%$ a day ---
both lose exactly $60$ gallons on day 1, which kills ``exponential just means faster,'' and the
linear model goes \emph{negative} by day 14 while the exponential one never reaches zero. That last
contrast is Lesson~6.0's asymptote, in context.
\end{multicols}
\end{skillbox}

\begin{skillbox}[Individual Work \& Assessment]{redbox}
\textbf{Exit Ticket} (independent, no notes), four items on one page: \textbf{(1)} the bare
translation --- write $b$ for ``grows $7\%$'' and for ``loses $30\%$''; \textbf{(2)} build and use a
model --- a \$$2400$ book collection gaining $7\%$ a year, giving $a$, $b$, $V(t)$, $V(3)\approx
\$2940.10$, \emph{and} a sentence interpreting it (this is the \textbf{A2.F.2f} evidence; a bare
number does not count); \textbf{(3)} build from a table ($5,20,80,320$), the item with no translation
at all; and \textbf{(4)} an \textbf{SOL-style multiple-choice} item --- a town of $4000$ decreasing
$6\%$ a year --- whose distractors are the three errors of the day: $4000(1.06)^{t}$ is the sign
error, $4000(0.06)^{t}$ is percent-as-base, and $4000-0.06t$ is the linear reflex. Sort the collected
tickets into \emph{got it / sign error / percent as base / linear reflex}. The two middle piles are
the same misconception at different strengths and are worth a three-minute re-teach off the Warm-Up
jacket. Anyone who misses item 3 needs the constant-ratio idea itself retaught, not the percent rule.
\end{skillbox}

\begin{skillbox}[Reinforcement \& Extension]{goldbox}
\textbf{Homework:} the translation in \emph{both} directions --- five sentences to convert into
factors, then three finished models to describe in words (the $0.85 \Rightarrow 15\%$
\emph{decrease} row is where the sign error surfaces, and $1.5 \Rightarrow 50\%$ catches
percent-as-base); three tables where one (constant difference $+3$) is deliberately \emph{not}
exponential and must be caught with Lesson~6.0's fingerprint test; a pre-drawn decay graph with
$(0,8)$ and $(3,1)$, where $b^{3}=\frac18$ tempts students to answer $\frac13$; two stories built,
evaluated, and interpreted (bacteria at $+40\%$ an hour, a phone at $-22\%$ a year); and an
\textbf{error-analysis} item carrying the lesson --- Jonah's $18{,}000(1.15)^{t}$ for a $15\%$
\emph{loss} and Priya's $250(0.07)^{w}$ for $7\%$ \emph{growth}, each to be named, described, and
corrected. \textbf{Extension:} two towns over ten years --- Ashford at $20{,}000$ losing $2\%$ a
year, Brookvale at $12{,}000$ gaining $5\%$ --- build both, verify the given table, locate the
crossing between years $6$ and $8$, and then \emph{fail} to solve
$20{,}000(0.98)^{t}=12{,}000(1.05)^{t}$ for lack of a common base. That failure is deliberate: it is
the wall Lesson~6.4 ends on, arriving a few days early, and it should be left standing. The last
part reruns Lesson~6.0's believability question on Ashford's asymptote. \textbf{Preview:}
Lesson~6.2, \emph{Graphing Exponential Functions \& Transformations} --- where $f(x)+k$ turns out to
be the one transformation that \emph{moves the asymptote}, and where today's
$8\left(\frac12\right)^{x}=8\cdot2^{-x}$ becomes the reflection rule.
\end{skillbox}
\end{document}
